% Generated draft: requires mathematical and source review.
\documentclass[11pt]{article}
\usepackage[margin=25mm]{geometry}
\usepackage{fontspec}
\setmainfont{Latin Modern Roman}
\usepackage{amsmath,amssymb,mathtools,unicode-math}
\setmathfont{Latin Modern Math}
\usepackage{enumitem,needspace,xurl,hyperref}
\setlength{\parindent}{0pt}
\setlength{\parskip}{5pt}
\setlength{\emergencystretch}{3em}
\newcommand{\N}{\mathbb{N}}
\newcommand{\Z}{\mathbb{Z}}
\newcommand{\Q}{\mathbb{Q}}
\newcommand{\R}{\mathbb{R}}
\newcommand{\C}{\mathbb{C}}
\newcommand{\F}{\mathbb{F}}
\newcommand{\A}{\mathbb{A}}
\newcommand{\PP}{\mathbb{P}}
\newcommand{\E}{\mathbb{E}}
\newcommand{\eps}{\varepsilon}
\newcommand{\dd}{\,\mathrm{d}}
\begin{document}
\section*{UnsolvedMath ID: 20003151}
\textbf{Source problem code:} catalogue entry
\subsection{Definitions and mathematical statement}
Fix a compact Lie group $G$ and a closed subgroup $H$. A metric on the fixed homogeneous presentation $G/H$ is $G$-invariant and is determined by an $\operatorname{Ad}(H)$-invariant positive-definite inner product on $\mathfrak g/\mathfrak h$. Problem 26 asks to classify those inner products whose resulting sectional curvatures are all nonnegative. Classification up to automorphism, equivariant isometry, arbitrary isometry, or scale should be distinguished; the original sentence does not choose an equivalence relation. Its terminal ``3'' is unattached extraction text and gives no stated dimension hypothesis.

The supplied partial classification concerns $SU(2)/\{e\}$ only. Fix a Lie-algebra basis $E_1,E_2,E_3$ with $[E_1,E_2]=2E_3$ and cyclic analogues. Up to a Lie-group automorphism, every left-invariant metric is diagonal,
\[
 \langle E_i,E_j\rangle=x_i\delta_{ij},\qquad
 0<x_1\leq x_2\leq x_3.
\]
In this bracket normalization the principal sectional curvature of $E_i\wedge E_j$, with remaining index $k$, has numerator
\[
 N_k=(x_i-x_j)^2+2x_k(x_i+x_j)-3x_k^2
\]
and positive denominator $x_1x_2x_3$. The curvature operator is diagonal in the corresponding principal two-plane basis. In the ordered chamber $N_1,N_2$ are strictly positive, so the aid's complete criterion for this presentation is
\[
 \sec\geq0\quad\Longleftrightarrow\quad
 (x_1-x_2)^2+2x_3(x_1+x_2)-3x_3^2\geq0,
\]
equivalently
\[
 x_3\leq\frac{x_1+x_2+2\sqrt{x_1^2-x_1x_2+x_2^2}}3.
\]
Strict inequality gives positive sectional curvature. Equality gives exactly one flat tangent two-plane at the identity, $\operatorname{span}(E_1,E_2)$, and its left translates; the other principal curvatures stay positive. The aid's sharp consequences are
\[
 x_3-x_2\leq x_1/3,\qquad x_3/x_2\leq4/3.
\]
The eigenvalue proportions $(1,1,4/3)$ attain both bounds. Rescaling all $x_i$ by a common positive factor does not affect the curvature-sign criterion.

This is a complete left-invariant $SU(2)$ answer up to automorphism, but does not classify nonnegatively curved homogeneous metrics on arbitrary compact presentations $G/H$. The record supplies no general classification beyond this case.
\subsection{Short English statement}
Classify the nonnegatively curved invariant metrics on a fixed compact homogeneous space. The supplied partial result treats left-invariant $SU(2)$ metrics: after ordering their three metric eigenvalues, one explicit inequality determines the full nonnegative-curvature region, its unique flat boundary plane and sharp spectral-gap bounds.
\subsection{Sources}
\begin{itemize}
\item Local UnsolvedMath registry, permanent ID $20003151$, source code catalogue entry; canonical location \texttt{aim-topology-notes.json}, \texttt{notes[238]}.
\item American Institute of Mathematics, \emph{Manifolds with nonnegative sectional curvature}, Problem 26, \url{https://aimath.org/WWN/nnsectcurvature/nnsectcurvature.pdf}. No proposer, remarks, or literature entries are supplied.
\item Attached research attempt 1 and summary refer to published three-inequality formulas for $SU(2)$ but give no author or publication reference. The attempt is marked partial result, proved, and low novelty confidence, with an explicit machine-generated, non-peer-reviewed research-aid warning requiring independent verification before citation.
\end{itemize}
\end{document}
